% ID: SEL-D0002-OB-E02
% ESTADO: cotejada. Fuente y transcripción contrastadas de forma independiente; salida revisada.
% Fuente: Selectividad/sel_2019_matematicas_aplicadas/Titular Septiembre/Exámenes Mat Apl CCSS II - A y B.pdf
% SHA-256: 31fb46947e141f5c8da8f7b381e7d27b67fe5a9d1597120454ce66fce0313f1a
% Página: 2; opción B; ejercicio 2.
% Transcripción de primera pasada basada en la imagen de la página original.
% La expresión dada es f'(x), no f(x).
% Sin solución ni modificación de datos o condiciones.

De una cierta función \(f\), sabemos que su función derivada es \(f'(x)=3x^2-3\).
\begin{enumerate}[label=\alph*),leftmargin=2.4em]
  \item \textbf{(1 punto)} Estudie los intervalos de crecimiento y decrecimiento de \(f\), y calcule la abscisa de sus extremos relativos.
  \item \textbf{(0.75 puntos)} Determine la curvatura de \(f\) y halle la abscisa de su punto de inflexión.
  \item \textbf{(0.75 puntos)} Calcule la función \(f\), sabiendo que su gráfica pasa por el punto \((-1,\,3)\).
\end{enumerate}
